% Options for packages loaded elsewhere
\PassOptionsToPackage{unicode}{hyperref}
\PassOptionsToPackage{hyphens}{url}
\documentclass[
]{article}
\usepackage{xcolor}
\usepackage[margin=1in]{geometry}
\usepackage{amsmath,amssymb}
\setcounter{secnumdepth}{-\maxdimen} % remove section numbering
\usepackage{iftex}
\ifPDFTeX
  \usepackage[T1]{fontenc}
  \usepackage[utf8]{inputenc}
  \usepackage{textcomp} % provide euro and other symbols
\else % if luatex or xetex
  \usepackage{unicode-math} % this also loads fontspec
  \defaultfontfeatures{Scale=MatchLowercase}
  \defaultfontfeatures[\rmfamily]{Ligatures=TeX,Scale=1}
\fi
\usepackage{lmodern}
\ifPDFTeX\else
  % xetex/luatex font selection
\fi
% Use upquote if available, for straight quotes in verbatim environments
\IfFileExists{upquote.sty}{\usepackage{upquote}}{}
\IfFileExists{microtype.sty}{% use microtype if available
  \usepackage[]{microtype}
  \UseMicrotypeSet[protrusion]{basicmath} % disable protrusion for tt fonts
}{}
\makeatletter
\@ifundefined{KOMAClassName}{% if non-KOMA class
  \IfFileExists{parskip.sty}{%
    \usepackage{parskip}
  }{% else
    \setlength{\parindent}{0pt}
    \setlength{\parskip}{6pt plus 2pt minus 1pt}}
}{% if KOMA class
  \KOMAoptions{parskip=half}}
\makeatother
\usepackage{color}
\usepackage{fancyvrb}
\newcommand{\VerbBar}{|}
\newcommand{\VERB}{\Verb[commandchars=\\\{\}]}
\DefineVerbatimEnvironment{Highlighting}{Verbatim}{commandchars=\\\{\}}
% Add ',fontsize=\small' for more characters per line
\usepackage{framed}
\definecolor{shadecolor}{RGB}{248,248,248}
\newenvironment{Shaded}{\begin{snugshade}}{\end{snugshade}}
\newcommand{\AlertTok}[1]{\textcolor[rgb]{0.94,0.16,0.16}{#1}}
\newcommand{\AnnotationTok}[1]{\textcolor[rgb]{0.56,0.35,0.01}{\textbf{\textit{#1}}}}
\newcommand{\AttributeTok}[1]{\textcolor[rgb]{0.13,0.29,0.53}{#1}}
\newcommand{\BaseNTok}[1]{\textcolor[rgb]{0.00,0.00,0.81}{#1}}
\newcommand{\BuiltInTok}[1]{#1}
\newcommand{\CharTok}[1]{\textcolor[rgb]{0.31,0.60,0.02}{#1}}
\newcommand{\CommentTok}[1]{\textcolor[rgb]{0.56,0.35,0.01}{\textit{#1}}}
\newcommand{\CommentVarTok}[1]{\textcolor[rgb]{0.56,0.35,0.01}{\textbf{\textit{#1}}}}
\newcommand{\ConstantTok}[1]{\textcolor[rgb]{0.56,0.35,0.01}{#1}}
\newcommand{\ControlFlowTok}[1]{\textcolor[rgb]{0.13,0.29,0.53}{\textbf{#1}}}
\newcommand{\DataTypeTok}[1]{\textcolor[rgb]{0.13,0.29,0.53}{#1}}
\newcommand{\DecValTok}[1]{\textcolor[rgb]{0.00,0.00,0.81}{#1}}
\newcommand{\DocumentationTok}[1]{\textcolor[rgb]{0.56,0.35,0.01}{\textbf{\textit{#1}}}}
\newcommand{\ErrorTok}[1]{\textcolor[rgb]{0.64,0.00,0.00}{\textbf{#1}}}
\newcommand{\ExtensionTok}[1]{#1}
\newcommand{\FloatTok}[1]{\textcolor[rgb]{0.00,0.00,0.81}{#1}}
\newcommand{\FunctionTok}[1]{\textcolor[rgb]{0.13,0.29,0.53}{\textbf{#1}}}
\newcommand{\ImportTok}[1]{#1}
\newcommand{\InformationTok}[1]{\textcolor[rgb]{0.56,0.35,0.01}{\textbf{\textit{#1}}}}
\newcommand{\KeywordTok}[1]{\textcolor[rgb]{0.13,0.29,0.53}{\textbf{#1}}}
\newcommand{\NormalTok}[1]{#1}
\newcommand{\OperatorTok}[1]{\textcolor[rgb]{0.81,0.36,0.00}{\textbf{#1}}}
\newcommand{\OtherTok}[1]{\textcolor[rgb]{0.56,0.35,0.01}{#1}}
\newcommand{\PreprocessorTok}[1]{\textcolor[rgb]{0.56,0.35,0.01}{\textit{#1}}}
\newcommand{\RegionMarkerTok}[1]{#1}
\newcommand{\SpecialCharTok}[1]{\textcolor[rgb]{0.81,0.36,0.00}{\textbf{#1}}}
\newcommand{\SpecialStringTok}[1]{\textcolor[rgb]{0.31,0.60,0.02}{#1}}
\newcommand{\StringTok}[1]{\textcolor[rgb]{0.31,0.60,0.02}{#1}}
\newcommand{\VariableTok}[1]{\textcolor[rgb]{0.00,0.00,0.00}{#1}}
\newcommand{\VerbatimStringTok}[1]{\textcolor[rgb]{0.31,0.60,0.02}{#1}}
\newcommand{\WarningTok}[1]{\textcolor[rgb]{0.56,0.35,0.01}{\textbf{\textit{#1}}}}
\usepackage{longtable,booktabs,array}
\newcounter{none} % for unnumbered tables
\usepackage{calc} % for calculating minipage widths
% Correct order of tables after \paragraph or \subparagraph
\usepackage{etoolbox}
\makeatletter
\patchcmd\longtable{\par}{\if@noskipsec\mbox{}\fi\par}{}{}
\makeatother
% Allow footnotes in longtable head/foot
\IfFileExists{footnotehyper.sty}{\usepackage{footnotehyper}}{\usepackage{footnote}}
\makesavenoteenv{longtable}
\usepackage{graphicx}
\makeatletter
\newsavebox\pandoc@box
\newcommand*\pandocbounded[1]{% scales image to fit in text height/width
  \sbox\pandoc@box{#1}%
  \Gscale@div\@tempa{\textheight}{\dimexpr\ht\pandoc@box+\dp\pandoc@box\relax}%
  \Gscale@div\@tempb{\linewidth}{\wd\pandoc@box}%
  \ifdim\@tempb\p@<\@tempa\p@\let\@tempa\@tempb\fi% select the smaller of both
  \ifdim\@tempa\p@<\p@\scalebox{\@tempa}{\usebox\pandoc@box}%
  \else\usebox{\pandoc@box}%
  \fi%
}
% Set default figure placement to htbp
\def\fps@figure{htbp}
\makeatother
\setlength{\emergencystretch}{3em} % prevent overfull lines
\providecommand{\tightlist}{%
  \setlength{\itemsep}{0pt}\setlength{\parskip}{0pt}}
\usepackage{bookmark}
\IfFileExists{xurl.sty}{\usepackage{xurl}}{} % add URL line breaks if available
\urlstyle{same}
\hypersetup{
  pdftitle={Simplex},
  pdfauthor={Prof.~Dr.~Stefan Böcker},
  hidelinks,
  pdfcreator={LaTeX via pandoc}}

\title{Simplex}
\author{Prof.~Dr.~Stefan Böcker}
\date{2025-11-28}

\begin{document}
\maketitle

\begin{Shaded}
\begin{Highlighting}[]
\DocumentationTok{\#\# Beispiel:}
\DocumentationTok{\#\# max z = 3 x1 + 2 x2}
\DocumentationTok{\#\# s.t.  x1 + x2 \textless{}= 4}
\DocumentationTok{\#\#      2x1 + x2 \textless{}= 5}
\DocumentationTok{\#\#      x1, x2 \textgreater{}= 0}

\CommentTok{\# Koeffizientenmatrix der Nebenbedingungen (inkl. Schlupfvariablen)}
\NormalTok{A }\OtherTok{\textless{}{-}} \FunctionTok{matrix}\NormalTok{(}\FunctionTok{c}\NormalTok{(}\DecValTok{1}\NormalTok{, }\DecValTok{1}\NormalTok{, }\DecValTok{1}\NormalTok{, }\DecValTok{0}\NormalTok{,}
              \DecValTok{2}\NormalTok{, }\DecValTok{1}\NormalTok{, }\DecValTok{0}\NormalTok{, }\DecValTok{1}\NormalTok{),}
            \AttributeTok{nrow =} \DecValTok{2}\NormalTok{, }\AttributeTok{byrow =} \ConstantTok{TRUE}\NormalTok{)}

\CommentTok{\# Rechte Seiten}
\NormalTok{b }\OtherTok{\textless{}{-}} \FunctionTok{c}\NormalTok{(}\DecValTok{4}\NormalTok{, }\DecValTok{5}\NormalTok{)}

\CommentTok{\# Zielfunktion: 3 x1 + 2 x2  (Slack{-}Koeffizienten = 0)}
\NormalTok{c\_vec }\OtherTok{\textless{}{-}} \FunctionTok{c}\NormalTok{(}\DecValTok{3}\NormalTok{, }\DecValTok{2}\NormalTok{, }\DecValTok{0}\NormalTok{, }\DecValTok{0}\NormalTok{)}

\CommentTok{\# Hilfsfunktion: Simplex{-}Tableau ausgeben}
\NormalTok{print\_tableau }\OtherTok{\textless{}{-}} \ControlFlowTok{function}\NormalTok{(A, b, c\_vec, }\AttributeTok{z0 =} \DecValTok{0}\NormalTok{, basis\_idx, }\AttributeTok{iter =} \DecValTok{0}\NormalTok{) \{}
\NormalTok{  m }\OtherTok{\textless{}{-}} \FunctionTok{nrow}\NormalTok{(A)}
\NormalTok{  n }\OtherTok{\textless{}{-}} \FunctionTok{ncol}\NormalTok{(A)}
  \FunctionTok{cat}\NormalTok{(}\StringTok{"===== Simplex{-}Tableau, Iteration"}\NormalTok{, iter, }\StringTok{"=====}\SpecialCharTok{\textbackslash{}n}\StringTok{"}\NormalTok{)}
\NormalTok{  tab }\OtherTok{\textless{}{-}} \FunctionTok{cbind}\NormalTok{(A, b)}
  \FunctionTok{colnames}\NormalTok{(tab) }\OtherTok{\textless{}{-}} \FunctionTok{c}\NormalTok{(}\FunctionTok{paste0}\NormalTok{(}\StringTok{"x"}\NormalTok{, }\DecValTok{1}\SpecialCharTok{:}\NormalTok{n), }\StringTok{"RHS"}\NormalTok{)}
  \FunctionTok{rownames}\NormalTok{(tab) }\OtherTok{\textless{}{-}} \FunctionTok{paste0}\NormalTok{(}\StringTok{"B"}\NormalTok{, }\DecValTok{1}\SpecialCharTok{:}\NormalTok{m)}
  \FunctionTok{print}\NormalTok{(}\FunctionTok{round}\NormalTok{(tab, }\DecValTok{3}\NormalTok{))}
\NormalTok{  z\_row }\OtherTok{\textless{}{-}} \FunctionTok{c}\NormalTok{(}\SpecialCharTok{{-}}\NormalTok{c\_vec, z0)  }\CommentTok{\# übliche Schul{-}Notation: Zielfunktion negiert}
  \FunctionTok{names}\NormalTok{(z\_row) }\OtherTok{\textless{}{-}} \FunctionTok{c}\NormalTok{(}\FunctionTok{paste0}\NormalTok{(}\StringTok{"x"}\NormalTok{, }\DecValTok{1}\SpecialCharTok{:}\NormalTok{n), }\StringTok{"RHS"}\NormalTok{)}
  \FunctionTok{cat}\NormalTok{(}\StringTok{"Zielfunktion (z{-}Zeile):}\SpecialCharTok{\textbackslash{}n}\StringTok{"}\NormalTok{)}
  \FunctionTok{print}\NormalTok{(}\FunctionTok{round}\NormalTok{(z\_row, }\DecValTok{3}\NormalTok{))}
  \FunctionTok{cat}\NormalTok{(}\StringTok{"Basisvariablen:"}\NormalTok{, }\FunctionTok{paste}\NormalTok{(}\FunctionTok{paste0}\NormalTok{(}\StringTok{"x"}\NormalTok{, basis\_idx), }\AttributeTok{collapse =} \StringTok{", "}\NormalTok{), }\StringTok{"}\SpecialCharTok{\textbackslash{}n\textbackslash{}n}\StringTok{"}\NormalTok{)}
\NormalTok{\}}


\NormalTok{simplex\_step }\OtherTok{\textless{}{-}} \ControlFlowTok{function}\NormalTok{(A, b, c\_vec, basis\_idx, }\AttributeTok{z0 =} \DecValTok{0}\NormalTok{, }\AttributeTok{iter =} \DecValTok{0}\NormalTok{) \{}
\NormalTok{  m }\OtherTok{\textless{}{-}} \FunctionTok{nrow}\NormalTok{(A)}
\NormalTok{  n }\OtherTok{\textless{}{-}} \FunctionTok{ncol}\NormalTok{(A)}

  \CommentTok{\# Tableau ausgeben}
  \FunctionTok{print\_tableau}\NormalTok{(A, b, c\_vec, z0, basis\_idx, iter)}

  \CommentTok{\# reduzierte Kosten (hier in Schul{-}Notation einfach {-}c)}
\NormalTok{  red\_costs }\OtherTok{\textless{}{-}} \SpecialCharTok{{-}}\NormalTok{c\_vec}

  \CommentTok{\# Optimalitätstest: keine negativen reduzierten Kosten in Nichtbasis}
\NormalTok{  nonbasis\_idx }\OtherTok{\textless{}{-}} \FunctionTok{setdiff}\NormalTok{(}\DecValTok{1}\SpecialCharTok{:}\NormalTok{n, basis\_idx)}
  \ControlFlowTok{if}\NormalTok{ (}\FunctionTok{all}\NormalTok{(red\_costs[nonbasis\_idx] }\SpecialCharTok{\textgreater{}=} \SpecialCharTok{{-}}\FloatTok{1e{-}9}\NormalTok{)) \{}
    \FunctionTok{cat}\NormalTok{(}\StringTok{"Optimalität erreicht.}\SpecialCharTok{\textbackslash{}n}\StringTok{"}\NormalTok{)}
    \FunctionTok{return}\NormalTok{(}\FunctionTok{list}\NormalTok{(}\AttributeTok{A =}\NormalTok{ A, }\AttributeTok{b =}\NormalTok{ b, }\AttributeTok{c =}\NormalTok{ c\_vec, }\AttributeTok{basis =}\NormalTok{ basis\_idx, }\AttributeTok{z0 =}\NormalTok{ z0,}
                \AttributeTok{optimal =} \ConstantTok{TRUE}\NormalTok{))}
\NormalTok{  \}}

  \CommentTok{\# Eintrittsvariable: Spalte mit minimalem reduzierten Kostenwert}
\NormalTok{  entering }\OtherTok{\textless{}{-}}\NormalTok{ nonbasis\_idx[}\FunctionTok{which.min}\NormalTok{(red\_costs[nonbasis\_idx])]}

  \CommentTok{\# Quotiententest}
\NormalTok{  col\_enter }\OtherTok{\textless{}{-}}\NormalTok{ A[, entering]}
\NormalTok{  valid }\OtherTok{\textless{}{-}}\NormalTok{ col\_enter }\SpecialCharTok{\textgreater{}} \FloatTok{1e{-}9}
  \ControlFlowTok{if}\NormalTok{ (}\SpecialCharTok{!}\FunctionTok{any}\NormalTok{(valid)) \{}
    \FunctionTok{stop}\NormalTok{(}\StringTok{"Problem unbeschränkt (keine gültige Pivotzeile)."}\NormalTok{)}
\NormalTok{  \}}
\NormalTok{  ratios }\OtherTok{\textless{}{-}}\NormalTok{ b[valid] }\SpecialCharTok{/}\NormalTok{ col\_enter[valid]}
\NormalTok{  pivot\_row }\OtherTok{\textless{}{-}} \FunctionTok{which}\NormalTok{(valid)[}\FunctionTok{which.min}\NormalTok{(ratios)]}

  \FunctionTok{cat}\NormalTok{(}\StringTok{"Eintrittsvariable: x"}\NormalTok{, entering,}
      \StringTok{"  (Spalte "}\NormalTok{, entering, }\StringTok{")}\SpecialCharTok{\textbackslash{}n}\StringTok{"}\NormalTok{, }\AttributeTok{sep =} \StringTok{""}\NormalTok{)}
  \FunctionTok{cat}\NormalTok{(}\StringTok{"Austrittsvariable: x"}\NormalTok{, basis\_idx[pivot\_row],}
      \StringTok{"  (Zeile "}\NormalTok{, pivot\_row, }\StringTok{")}\SpecialCharTok{\textbackslash{}n\textbackslash{}n}\StringTok{"}\NormalTok{, }\AttributeTok{sep =} \StringTok{""}\NormalTok{)}

  \CommentTok{\# Pivot{-}Operation}
\NormalTok{  pivot }\OtherTok{\textless{}{-}}\NormalTok{ A[pivot\_row, entering]}
\NormalTok{  A[pivot\_row, ] }\OtherTok{\textless{}{-}}\NormalTok{ A[pivot\_row, ] }\SpecialCharTok{/}\NormalTok{ pivot}
\NormalTok{  b[pivot\_row]   }\OtherTok{\textless{}{-}}\NormalTok{ b[pivot\_row]   }\SpecialCharTok{/}\NormalTok{ pivot}

  \ControlFlowTok{for}\NormalTok{ (i }\ControlFlowTok{in} \FunctionTok{seq\_len}\NormalTok{(m)) \{}
    \ControlFlowTok{if}\NormalTok{ (i }\SpecialCharTok{==}\NormalTok{ pivot\_row) }\ControlFlowTok{next}
\NormalTok{    factor }\OtherTok{\textless{}{-}}\NormalTok{ A[i, entering]}
\NormalTok{    A[i, ] }\OtherTok{\textless{}{-}}\NormalTok{ A[i, ] }\SpecialCharTok{{-}}\NormalTok{ factor }\SpecialCharTok{*}\NormalTok{ A[pivot\_row, ]}
\NormalTok{    b[i]   }\OtherTok{\textless{}{-}}\NormalTok{ b[i]   }\SpecialCharTok{{-}}\NormalTok{ factor }\SpecialCharTok{*}\NormalTok{ b[pivot\_row]}
\NormalTok{  \}}

  \CommentTok{\# Zielfunktion anpassen (einfache Variante)}
\NormalTok{  factor\_z }\OtherTok{\textless{}{-}}\NormalTok{ c\_vec[entering]}
\NormalTok{  c\_vec }\OtherTok{\textless{}{-}}\NormalTok{ c\_vec }\SpecialCharTok{{-}}\NormalTok{ factor\_z }\SpecialCharTok{*}\NormalTok{ A[pivot\_row, ]}
\NormalTok{  z0    }\OtherTok{\textless{}{-}}\NormalTok{ z0 }\SpecialCharTok{+}\NormalTok{ factor\_z }\SpecialCharTok{*}\NormalTok{ b[pivot\_row]}

  \CommentTok{\# Basis aktualisieren}
\NormalTok{  basis\_idx[pivot\_row] }\OtherTok{\textless{}{-}}\NormalTok{ entering}

  \FunctionTok{list}\NormalTok{(}\AttributeTok{A =}\NormalTok{ A, }\AttributeTok{b =}\NormalTok{ b, }\AttributeTok{c =}\NormalTok{ c\_vec, }\AttributeTok{basis =}\NormalTok{ basis\_idx, }\AttributeTok{z0 =}\NormalTok{ z0,}
       \AttributeTok{optimal =} \ConstantTok{FALSE}\NormalTok{)}
\NormalTok{\}}


\CommentTok{\# Startbasis: Schlupfvariablen x3, x4}
\NormalTok{basis\_idx }\OtherTok{\textless{}{-}} \FunctionTok{c}\NormalTok{(}\DecValTok{3}\NormalTok{, }\DecValTok{4}\NormalTok{)}

\NormalTok{state }\OtherTok{\textless{}{-}} \FunctionTok{list}\NormalTok{(}\AttributeTok{A =}\NormalTok{ A, }\AttributeTok{b =}\NormalTok{ b, }\AttributeTok{c =}\NormalTok{ c\_vec,}
              \AttributeTok{basis =}\NormalTok{ basis\_idx, }\AttributeTok{z0 =} \DecValTok{0}\NormalTok{)}

\NormalTok{iter }\OtherTok{\textless{}{-}} \DecValTok{0}
\ControlFlowTok{repeat}\NormalTok{ \{}
\NormalTok{  res }\OtherTok{\textless{}{-}} \FunctionTok{simplex\_step}\NormalTok{(state}\SpecialCharTok{$}\NormalTok{A, state}\SpecialCharTok{$}\NormalTok{b, state}\SpecialCharTok{$}\NormalTok{c,}
\NormalTok{                      state}\SpecialCharTok{$}\NormalTok{basis, state}\SpecialCharTok{$}\NormalTok{z0, }\AttributeTok{iter =}\NormalTok{ iter)}
  \ControlFlowTok{if}\NormalTok{ (}\FunctionTok{isTRUE}\NormalTok{(res}\SpecialCharTok{$}\NormalTok{optimal)) \{}
\NormalTok{    state }\OtherTok{\textless{}{-}}\NormalTok{ res}
    \ControlFlowTok{break}
\NormalTok{  \}}
\NormalTok{  state }\OtherTok{\textless{}{-}}\NormalTok{ res}
\NormalTok{  iter }\OtherTok{\textless{}{-}}\NormalTok{ iter }\SpecialCharTok{+} \DecValTok{1}
\NormalTok{\}}
\end{Highlighting}
\end{Shaded}

\begin{verbatim}
## ===== Simplex-Tableau, Iteration 0 =====
##    x1 x2 x3 x4 RHS
## B1  1  1  1  0   4
## B2  2  1  0  1   5
## Zielfunktion (z-Zeile):
##  x1  x2  x3  x4 RHS 
##  -3  -2   0   0   0 
## Basisvariablen: x3, x4 
## 
## Eintrittsvariable: x1  (Spalte 1)
## Austrittsvariable: x4  (Zeile 2)
## 
## ===== Simplex-Tableau, Iteration 1 =====
##    x1  x2 x3   x4 RHS
## B1  0 0.5  1 -0.5 1.5
## B2  1 0.5  0  0.5 2.5
## Zielfunktion (z-Zeile):
##   x1   x2   x3   x4  RHS 
##  0.0 -0.5  0.0  1.5  7.5 
## Basisvariablen: x3, x1 
## 
## Eintrittsvariable: x2  (Spalte 2)
## Austrittsvariable: x3  (Zeile 1)
## 
## ===== Simplex-Tableau, Iteration 2 =====
##    x1 x2 x3 x4 RHS
## B1  0  1  2 -1   3
## B2  1  0 -1  1   1
## Zielfunktion (z-Zeile):
##  x1  x2  x3  x4 RHS 
##   0   0   1   1   9 
## Basisvariablen: x2, x1 
## 
## Optimalität erreicht.
\end{verbatim}

\begin{Shaded}
\begin{Highlighting}[]
\CommentTok{\# Optimale Lösung aus dem finalen Tableau ablesen}
\NormalTok{m }\OtherTok{\textless{}{-}} \FunctionTok{nrow}\NormalTok{(state}\SpecialCharTok{$}\NormalTok{A); n }\OtherTok{\textless{}{-}} \FunctionTok{ncol}\NormalTok{(state}\SpecialCharTok{$}\NormalTok{A)}
\NormalTok{x }\OtherTok{\textless{}{-}} \FunctionTok{rep}\NormalTok{(}\DecValTok{0}\NormalTok{, n)}
\ControlFlowTok{for}\NormalTok{ (i }\ControlFlowTok{in} \DecValTok{1}\SpecialCharTok{:}\NormalTok{m) \{}
\NormalTok{  x[state}\SpecialCharTok{$}\NormalTok{basis[i]] }\OtherTok{\textless{}{-}}\NormalTok{ state}\SpecialCharTok{$}\NormalTok{b[i]}
\NormalTok{\}}
\FunctionTok{cat}\NormalTok{(}\StringTok{"Optimale Variablebelegung:}\SpecialCharTok{\textbackslash{}n}\StringTok{"}\NormalTok{)}
\end{Highlighting}
\end{Shaded}

\begin{verbatim}
## Optimale Variablebelegung:
\end{verbatim}

\begin{Shaded}
\begin{Highlighting}[]
\FunctionTok{print}\NormalTok{(}\FunctionTok{round}\NormalTok{(x, }\DecValTok{3}\NormalTok{))}
\end{Highlighting}
\end{Shaded}

\begin{verbatim}
## [1] 1 3 0 0
\end{verbatim}

\begin{Shaded}
\begin{Highlighting}[]
\FunctionTok{cat}\NormalTok{(}\StringTok{"Optimaler Zielfunktionswert z* ="}\NormalTok{, }\FunctionTok{round}\NormalTok{(state}\SpecialCharTok{$}\NormalTok{z0, }\DecValTok{3}\NormalTok{), }\StringTok{"}\SpecialCharTok{\textbackslash{}n}\StringTok{"}\NormalTok{)}
\end{Highlighting}
\end{Shaded}

\begin{verbatim}
## Optimaler Zielfunktionswert z* = 9
\end{verbatim}

{\def\LTcaptype{none} % do not increment counter
\begin{longtable}[]{@{}llllll@{}}
\toprule\noalign{}
& x₁ & x₂ & s₁ & s₂ & RHS \\
\midrule\noalign{}
\endhead
\bottomrule\noalign{}
\endlastfoot
s₁ & 1 & 1 & 1 & 0 & 4 \\
s₂ & 2 & 1 & 0 & 1 & 5 \\
z & -3 & -2 & 0 & 0 & 0 \\
\end{longtable}
}

\end{document}
